\documentclass[../../main.tex]{subfiles}% 注意这里的文件路径不能用 ./main.tex ,否则用latexmk编译子文件会报错
\graphicspath{{\subfix{./image/}}{\subfix{../../image/}}} % 指定图片引用路径,后续可以直接使用图片文件名
% 如果图片存放在上级目录,则使用 ../../image/ 作为备用路径

% 例如:
% \begin{figure}[H]
% \centering
% \includegraphics[width=0.3\textwidth]{图.png}
% \caption{}
% \label{figure:图}
% \end{figure}
% 注意:上述\label{}一定要放在\caption{}之后,否则引用图片序号会只会显示??.

\begin{document}

\section{幂级数的全纯开拓}

\begin{definition}
设 \(f\) 是域 \(D\) 上的全纯函数, 对于 \(\zeta \in \partial D\), 如果存在 \(\zeta\) 的邻域 \(B(\zeta,\delta)\) 及其上的全纯函数 \(g\), 使得 \(f(z)=g(z)\) 对 \(z \in D \cap B(\zeta,\delta)\) 成立, 就称 \(\zeta\) 是 \(f\) 的\textbf{正则点}, 否则称 \(\zeta\) 为 \(f\) 的\textbf{奇点}.
\end{definition}

\begin{examples}
幂级数 \(f(z)=\sum\limits_{n=0}^{\infty} z^n\) 的收敛半径 \(R=1\). 取 \(z_0 \in B(0,1), z_0 \neq 0\), 则 \(f\) 在 \(z_0\) 处的幂级数展开式为
\begin{align*}
\frac{1}{1-z} = \frac{1}{1-z_0} \frac{1}{1-\frac{z-z_0}{1-z_0}} = \sum_{n=0}^{\infty} \frac{1}{(1-z_0)^{n+1}}(z-z_0)^n.
\end{align*}
右端幂级数的收敛半径
\begin{align*}
R = \left[ \varlimsup\limits_{n \to \infty} \sqrt[n]{\frac{1}{|1-z_0|^{n+1}}} \right]^{-1} 
= |1-z_0| 
\geqslant 1-|z_0|,
\end{align*}
等号成立当且仅当 \(z_0\) 是非负实数. 这说明除了 \(z=1\) 外, \(f\) 都能通过其收敛圆周上的点全纯开拓出去. \(\square\)
\end{examples}

\begin{theorem}
幂级数的收敛圆周上必有其奇点.
\end{theorem}
\begin{proof}
设幂级数
\[
f(z)=\sum_{n=0}^{\infty} a_n z^n
\]
的收敛半径为 \(R\) (\(0<R<\infty\)), 那么 \(f\) 是 \(B(0,R)\) 中的全纯函数. 如果 \(\partial B(0,R)\) 上没有 \(f\) 的奇点, 那么对每个 \(\zeta \in \partial B(0,R)\), 存在 \(B(\zeta,r_\zeta)\) 及 \(B(\zeta,r_\zeta)\) 上的全纯函数 \(f_\zeta\), 使得当 \(z \in B(0,R) \cap B(\zeta,r_\zeta)\) 时, \(f(z)=f_\zeta(z)\). 显然, \(\{B(\zeta,r_\zeta)\}\) 的全体构成 \(\partial B(0,R)\) 的一个开覆盖, 由于 \(\partial B(0,R)\) 是紧集, 故从中可以选出有限个 \(B(\zeta_1,r_{\zeta_1}),\cdots,B(\zeta_m,r_{\zeta_m})\),
使得
\[
\partial B(0,R) \subseteq \bigcup\limits_{n=1}^{m} B(\zeta_n,r_{\zeta_n}).
\]
记
\[
D = B(0,R) \cup \left( \bigcup\limits_{n=1}^{m} B(\zeta_n,r_{\zeta_n}) \right),
\]
并在 \(D\) 上定义函数
\[
g(z) = \begin{cases} f(z), & z \in B(0,R); \\ f_{\zeta_k}(z), & z \in B(\zeta_k,r_{\zeta_k}), k=1,\cdots,m, \end{cases}
\]
则 \(g\) 是 \(D\) 上的单值全纯函数. 因若 \(B(\zeta_j,r_{\zeta_j}) \cap B(\zeta_k,r_{\zeta_k}) \neq \varnothing\), 则当 \(z \in B(\zeta_j,r_{\zeta_j}) \cap B(\zeta_k,r_{\zeta_k}) \cap B(0,R)\) 时, \(f_{\zeta_j}(z) = f_{\zeta_k}(z)\). 由唯一性定理, \(f_{\zeta_j}\) 与 \(f_{\zeta_k}\) 在 \(B(\zeta_j,r_{\zeta_j}) \cap B(\zeta_k,r_{\zeta_k})\) 上也相等. 所以, \(g\) 是 \(f\) 在 \(D\) 上的全纯开拓. 由于
\begin{align*}
g(z) = \sum_{n=0}^{\infty} \frac{g^{(n)}(0)}{n!} z^n = \sum_{n=0}^{\infty} \frac{f^{(n)}(0)}{n!} z^n = \sum_{n=0}^{\infty} a_n z^n,
\end{align*}
但它的收敛半径严格大于 \(R\), 这就与假定相矛盾.
\end{proof}

\begin{proposition}\label{proposition:幂级数收敛圆上必有奇点}
设幂级数
\[
f(z) = \sum_{n=0}^{\infty} a_n z^n
\]
的收敛半径为 \(R\) (\(0<R<\infty\)), 如果存在自然数 \(n_0\), 使得当 \(n \geqslant n_0\) 时总有 \(a_n \geqslant 0\), 那么 \(z=R\) 必是 \(f\) 的奇点.
\end{proposition}
\begin{proof}
如果 \(z=R\) 不是 \(f\) 的奇点, 那么 \(f\) 在 \(\frac{R}{2}\) 处展开的幂级数
\(\sum_{n=0}^{\infty} \frac{f^{(n)}\left( \frac{R}{2} \right)}{n!} \left( z-\frac{R}{2} \right)^n\) 的收敛半径 \(\rho > \frac{R}{2}\), 即
\begin{align}
\varlimsup\limits_{n \to \infty} \sqrt[n]{\left| \frac{f^{(n)}\left( \frac{R}{2} \right)}{n!} \right|} = \frac{1}{\rho} < \frac{2}{R}. \label{eq:A1b2C3d4}
\end{align}
由\cref{theorem:幂级数在其收敛圆内确定一个全纯函数并且任意阶可导}知
\begin{align*}
f^{(k)}\left( \frac{R}{2} \right) &= \sum_{n=k}^{\infty} n(n-1)\cdots(n-k+1)a_n \left( \frac{R}{2} \right)^{n-k}.
\end{align*}
由于当 \(n \geqslant n_0\) 时 \(a_n \geqslant 0\), 故当 \(k \geqslant n_0\) 时, 有
\begin{align*}
\left| f^{(k)}\left( \frac{R}{2} \right) \right| = \left| \sum_{n=k}^{\infty} n(n-1)\cdots(n-k+1)a_n \left( \frac{R}{2} \right)^{n-k} \right| 
= \sum_{n=k}^{\infty} n(n-1)\cdots(n-k+1)a_n \left( \frac{R}{2} \right)^{n-k}.
\end{align*}
今在收敛圆周 \(\partial B(0,R)\) 上任取一点 \(e^{\mathrm{i}\theta}R\), 则当 \(k \geqslant n_0\) 时, 有
\begin{align}
\left| f^{(k)}\left( \frac{1}{2} e^{\mathrm{i}\theta}R \right) \right| &= \left| \sum_{n=k}^{\infty} n(n-1)\cdots(n-k+1)a_n \left( \frac{e^{\mathrm{i}\theta}R}{2} \right)^{n-k} \right| \nonumber \\
&\leqslant \sum_{n=k}^{\infty} n(n-1)\cdots(n-k+1)a_n \left( \frac{R}{2} \right)^{n-k} \nonumber \\
&= \left| f^{(k)}\left( \frac{R}{2} \right) \right|. \label{eq:E5f6G7h8}
\end{align}
于是由 \eqref{eq:A1b2C3d4} 式和 \eqref{eq:E5f6G7h8} 式知道, \(f\) 在 \(\frac{1}{2} e^{\mathrm{i}\theta}R\) 处展开的幂级数的收敛半径为
\begin{align*}
\rho' = \left[ \varlimsup\limits_{n \to \infty} \sqrt[n]{\left| \frac{f^{(n)}\left( \frac{1}{2} e^{\mathrm{i}\theta}R \right)}{n!} \right|} \right]^{-1} \geqslant \left[ \varlimsup\limits_{n \to \infty} \sqrt[n]{\left| \frac{f^{(n)}\left( \frac{R}{2} \right)}{n!} \right|} \right]^{-1} > \frac{R}{2},
\end{align*}
这说明 \(e^{\mathrm{i}\theta}R\) 是 \(f\) 的正则点. 由于 \(e^{\mathrm{i}\theta}R\) 是 \(\partial B(0,R)\) 上的任意点, 所以收敛圆周 \(\partial B(0,R)\) 上没有 \(f\) 的奇点, 这与定理 6.2.3 的结论相矛盾！\end{proof}

\begin{examples}
幂级数 \(f(z)=\sum\limits_{n=0}^{\infty} z^{n!}\) 的收敛半径显然为 1, 我们证明收敛圆周 \(\partial B(0,1)\) 上的每一点都是它的奇点. 这时, 我们称 \(\partial B(0,1)\) 是 \(f\) 的\textbf{自然边界}.
\end{examples}
\begin{remark}
这个例子说明:幂级数在其收敛圆周上必有奇点, 但是不一定有正则点.
\end{remark}
\begin{proof}
由\cref{proposition:幂级数收敛圆上必有奇点}知道, \(z=1\) 是 \(f\) 的奇点. 今取既约分数 \(\frac{p}{q}\) \((q>0)\), 令 \(g(z) = f(e^{\mathrm{i}2\pi \frac{p}{q}} z)\), 则
\begin{align*}
g(z) = \sum_{n=0}^{\infty} e^{\mathrm{i}2\pi \frac{p}{q} n!} z^{n!} = \sum_{n=0}^{q-1} e^{\mathrm{i}2\pi \frac{p}{q} n!} z^{n!} + \sum_{n=q}^{\infty} z^{n!}.
\end{align*}
仍由\cref{proposition:幂级数收敛圆上必有奇点}知道, \(z=1\) 是 \(g\) 的奇点, 因而 \(z=e^{\mathrm{i}2\pi \frac{p}{q}}\) 是 \(f\) 的奇点. 现在证明 \(\partial B(0,1)\) 上的每一点都是 \(f\) 的奇点. 如果 \(\zeta \in \partial B(0,1)\) 不是 \(f\) 的奇点, 则必存在圆盘 \(B(\zeta,r)\), 使得 \(f\) 能全纯开拓到 \(B(\zeta,r)\), 故 \(\partial B(0,1) \cap B(\zeta,r)\) 中的每个点都是 \(f\) 的正则点. 而 \(\partial B(0,1) \cap B(\zeta,r)\) 中必有形如 \(e^{\mathrm{i}2\pi \frac{p}{q}}\) 的点, 这就和上面的结论相矛盾！\end{proof}

\begin{examples}
研究幂级数 \(f(z) = \sum\limits_{n=1}^{\infty} \frac{z^{n!}}{n^2}\) 的奇点.
\end{examples}
\begin{solution}
上述幂级数的收敛半径显然为 1, 并且在其闭收敛圆盘 \(\overline{B(0,1)}\) 上绝对一致收敛. 用与例 6.2.5 完全相同的方法, 知其收敛圆周上的每一点都是它的奇点, 即单位圆周为其自然边界. 
\end{solution}





\end{document}